\subsection{Проверка осуществимости авторской реализации SCC}
\label{sec:scc-feasibility}
Отдельный пилот разработки исполняет исходные классы ролей и контроллер Self-collaboration из авторской реализации 2024 года с адаптацией вызовов к Codex по подписке, GPT-5.6 Luna и уровню рассуждения medium. Три ранее использованные задачи разработки (/325, /322 и /1036), выбранные по их существующему порядку, дали три полные программы за 11 модельных вызовов. Все три программы проверены штатными тестами BigCodeBench: /325 прошла, /322 и /1036 не прошли. В той же фиксированной среде прошли все три эталона и были отклонены все три неправильных контроля. Известный расход составляет 48\,023 входных и выходных токена без пропущенных счётчиков и USD 0{,}0138646 условной API-оценки. Это не счёт подписки и не полная стоимость оценивания.

Для Luna используются собственные зарегистрированные базовые тарифы: USD 0{,}20/0{,}02/1{,}20 за миллион входных/кешированных/выходных токенов~\cite{openai56luna2026}:
\[
V_{\mathrm{Luna}}=\frac{0.20(I-C)+0.02C+1.20O}{10^6}.
\]
Здесь $C$ входит в $I$, а токены рассуждения --- в $O$. Коэффициенты отличаются от предыдущей серии с GPT-5.4 mini. Формула задаёт оценку по базовому тарифу; надбавки за длинные запросы и запись кеша требуют отдельной проверки применимости.

Статья 2023 года описывает тестировщика, который имитирует проверку и возвращает отчёт на естественном языке~\cite{dong2023selfcollaboration}. Авторский код 2024 года исполняет примеры, созданные моделью, и использует отсутствие исключения как внутренний сигнал остановки; напечатанные значения автоматически не сравниваются с ожидаемыми. Объектом воспроизведения служит именно эта версия кода~\cite{scc2024source}. Логика контроллера, запросы и извлечение кода сохранены, но предел составляет две итерации кодировщика вместо четырёх, заданных по умолчанию в исходном конструкторе: первоначальная реализация и не более одного исправления, до четырёх модельных вызовов с учётом аналитика и тестировщика. Сравнение ограничено этой конфигурацией. Сгенерированный код исполняется в изолированном контейнере. Адаптер сериализует сообщения ролей через CLI и не воспроизводит исходные параметры API-сэмплирования; полная задача передаётся текстом, ожидается самодостаточная программа. Повторная обработка исходных ответов восстанавливает переходы авторского контроллера и все три итоговые программы. В этом первоначальном пилоте выполнен один повтор без одновременного запуска методов сравнения.

В последующей проверочной серии 11 сентября 2026 года на тех же трёх ранее использованных задачах выполнены новые назначения SR, SN и SCC через общий транспорт Luna. Все девять программ проверены штатными тестами: SR решила 0/3 задач, SN --- 1/3, SCC --- 1/3. Сохранены ответы и счётчики расхода всех 18 модельных вызовов; их известная условная API-оценка составляет USD 0{,}02301176. Контрольные записи подтверждают успешные проверки эталонов и отказы неправильных решений в неизменной среде; итоговые программы SCC соответствуют переходам контроллера, записанным в этом пилоте. Эта серия проверяет работоспособность процедуры сравнения, включая обработку неудач; три ранее раскрытые задачи и один повтор не позволяют сделать статистический вывод о различиях методов. Задачи обоих пилотов исключены из выборки на 1\,000 задач.

\subsection{Завершённое сравнение SCC на 2\,000 блоках}
Дополнительная серия включает 6\,000 назначений методов на 956 задачах; 941 задача проходит контрольную проверку. Получены 1\,600 завершённых процедур SCC, 1\,997 SR и 1\,998 SN. Все 5\,595 завершённых программ проверены штатными тестами. Остальные 405 назначений остаются незавершёнными: у SCC сохранены 396 ошибок Docker и четыре приостановленные попытки, у SR --- три приостановленные попытки, у SN --- одна приостановленная попытка и одна ошибка потока ответа. Их качество остаётся неизвестным.

Для каждого сравнения сначала усредняются разности по совпадающим наблюдаемым повторам внутри задачи, затем все участвующие задачи получают одинаковый вес. Разность SCC минус SR составляет $-2{,}68$ процентного пункта по 883 задачам (95\%-ный интервал $[-4{,}72,-0{,}60]$). Для SCC минус SN она составляет $-2{,}48$ пункта по 882 задачам (интервал $[-4{,}50,-0{,}43]$). После поправки Холма в отдельном семействе из двух тестов оба значения $p$ равны $0{,}0217$. Эти результаты относятся к наблюдаемым допустимым парам. Предельные оценки с учётом неизвестных исходов на всех 941 допустимой задаче равны $[-12{,}98,+6{,}61]$ пункта относительно SR и $[-13{,}14,+6{,}43]$ относительно SN. Оба диапазона допускают любой знак; порядок качества во всей назначенной выборке не установлен.

Денежная оценка по тарифам API сравнивается для завершённых процедур с известными счётчиками. Отношение средних по парным задачам равно $3{,}138$ относительно SR (896 задач; 95\%-ный интервал $[3{,}067,3{,}208]$) и $2{,}939$ относительно SN (895 задач; интервал $[2{,}871,3{,}010]$). Это отношение двух выборочных средних, а не среднее индивидуальных отношений. Такая оценка не доказывает сохранение качества и не является счётом за подписку.

Отдельная диагностическая проверка охватывает все 399 сохранённых программ разработчика из незавершённых процессов SCC. Контроллер не перезапускается, сгенерированные проверки к программам не добавляются. Штатное оценивание даёт 196 успешных исходов, 201 неуспех и два превышения времени. После проверки эталонов качество определено для 392 программ: 196 проходят тесты и 196 не проходят; ещё семь исключены из оценки качества из-за контрольной проверки среды. Промежуточные программы не подменяют неизвестные исходы полных процессов SCC.
